% !TeX program = pdflatex
% =====================================================================
% ПЕРЕВОД НА АНГЛИЙСКИЙ ЯЗЫК --- отдельный файл по правилам журнала
% («Требования к рукописи после принятия к публикации»: перевод сведений
% первой страницы, списка источников и подписей к рисункам подается
% ОТДЕЛЬНО от рукописи; в main.tex по решению авторов сохранен английский
% титульный лист с аннотацией).
% Компилируется: pdflatex english_translation.tex (2 прохода)
% =====================================================================
\documentclass[11pt,twoside,a4paper]{article}
\usepackage{styleAND_2023}
\usepackage{xurl}% переносимые DOI
\emergencystretch=2em

\newcommand{\yearpub}{2026}
\newcommand{\volumemag}{xx}
\newcommand{\numbermag}{x}
\newcommand{\published}{xx.xx.xxxx}

\renewcommand{\papernumber}{xxxxxx}
\renewcommand{\EDN}{EDN:~XXXXXX}
\renewcommand{\received}{xx.xx.xxxx}
\renewcommand{\accepted}{xx.xx.xxxx}
\renewcommand{\onlinepublished}{xx.xx.xxxx}
\renewcommand{\firstpage}{xxx}
\renewcommand{\lastpage}{xxx}

% (счётчики теорем здесь не объявлены --- zerocouter не нужен)

\renewcommand{\titpaperen}{The effect of hemodynamic response nonlinearity on the accuracy of activation contrast estimation in fMRI}
\renewcommand{\authorsen}{Korelov~M.\,S., Knyazeva~I.\,S., Masharipov~R.\,S., Medvedeva~T.\,V., Kireev~M.\,V.}

\renewcommand{\authors}{\authorsen}

\begin{document}
\label{firstpage}
\thispagestyle{firststyle}
\vspace*{.2cm}
\infojournal
\vspace*{.8cm}

{\small \noindent
	Article\hfill \doi\\

}\vspace{.4cm}

\begin{center}\textbf{\large \titpaperen}
	\vspace{1\baselineskip}
	
\emph{M.\,S.~Korelov$^{1,3}$\lett, I.\,S.~Knyazeva$^{1}$, R.\,S.~Masharipov$^{1}$, T.\,V.~Medvedeva$^{1}$, M.\,V.~Kireev$^{1,2}$}\vspace{.7\baselineskip}

\footnotesize{
	$^{1}$Bekhtereva Institute of the Human Brain Russian Academy of Sciences, Russia \\
	$^{2}$St.Petersburg State University, Russia\\
	$^{3}$Peter the Great St.Petersburg Polytechnic University, Russia\\
	E-mail: \Letter\,korelovm@yandex.ru, knyazeva@ihb.spb.ru, masharipov@ihb.spb.ru, t.medvedeva2919@gmail.com, kireev@ihb.spb.ru\\[1pt]
	\emph{Received  \received, accepted \accepted, available online \onlinepublished, published \published}	
		
	}
\end{center}\vspace{.5\baselineskip}

\begin{abstract}

\noindent\textbf{\textit{Abstract.}}
Activation of brain regions is estimated from the fMRI BOLD signal by a general linear model with the canonical hemodynamic response function (HRF); condition differences are expressed as contrasts. Yet BOLD formation is nonlinear, beyond convolution with a fixed function.
The \textit{purpose} is to determine their accuracy and test one-step analytical corrections without refitting a nonlinear model.
\textit{Methods.} Synthetic BOLD signals are generated by P-DCM and Balloon--Windkessel models with a specified input, so true amplitudes are known. Each model's linear approximation separates HRF-shape error from the nonlinear contribution. Full and static corrections use the second-order Volterra expansion, fixed before testing on independent data.
\textit{Results.} Errors depend on event density and sequence and on overall stimulation level. At the largest amplitudes, median contrast underestimation is about 22\%. At zero contrast, differences in preceding events can produce a nonzero estimate. HRF-shape error and the nonlinear contribution can partly cancel. The full correction reduces the contrast RMSE by 40--80\% for clearly different amplitudes, depending on noise, and by 23--31\% under parameter spread; at zero contrast with random assignment the gain is small and strong noise raises the error; the static correction improves coefficients more often than the contrast. The second-order approximation reproduces the contrast error sign wherever resolvable.
\textit{Conclusion.} Standard analysis with the canonical HRF gives biased estimates depending on amplitudes and temporal structure of stimulation. Second-order corrections markedly reduce the error for large amplitude differences, but not in all conditions. The findings apply to the studied models and protocols; extension to humans requires separate justification.
\vspace{0.5\baselineskip}


\keywordsen{functional MRI, BOLD signal, hemodynamic response function, Balloon--Wind\-kes\-sel model, nonlinearity, general linear model, activation contrast, Volterra series.}

\Acknowledgementsen{TO BE FILLED: funding and acknowledgements.}

\Referencen{\textit{\authorsen}\hspace{3pt}\titpaperen.\hspace{2pt}{Izvestiya VUZ. Applied Nonlinear Dynamics}. \yearpub;\volumemag(\numbermag):\firstpage--\lastpage. \doi}

\licensen


\end{abstract}

{\sloppy
\begin{ltrtq}

	\bibitem{haxby2001}
	{Haxby~JV, Gobbini~MI, Furey~ML, Ishai~A, Schouten~JL, Pietrini~P.} Distributed and overlapping representations of faces and objects in ventral temporal cortex. {Science}. 2001;293(5539):2425--2430. \href{https://doi.org/10.1126/science.1063736}{DOI:~10.1126/science.1063736}.

	\bibitem{ogawa1990}
	{Ogawa~S, Lee~TM, Kay~AR, Tank~DW.} Brain magnetic resonance imaging with contrast dependent on blood oxygenation. {Proc. Natl. Acad. Sci. USA}. 1990;87(24):9868--9872. \href{https://doi.org/10.1073/pnas.87.24.9868}{DOI:~10.1073/\allowbreak pnas.87.24.9868}.

	\bibitem{kwong1992}
	{Kwong~KK, Belliveau~JW, Chesler~DA, et al.} Dynamic magnetic resonance imaging of human brain activity during primary sensory stimulation. {Proc. Natl. Acad. Sci. USA}. 1992;89(12):5675--5679. \href{https://doi.org/10.1073/pnas.89.12.5675}{DOI:~10.1073/\allowbreak pnas.89.12.5675}.

	\bibitem{friston2000}
	{Friston~KJ, Mechelli~A, Turner~R, Price~CJ.} Nonlinear responses in fMRI: the Balloon model, Volterra kernels, and other hemodynamics. {NeuroImage}. 2000;12(4):466--477. \href{https://doi.org/10.1006/nimg.2000.0630}{DOI:~10.1006/\allowbreak nimg.2000.0630}.

	\bibitem{buxton2004}
	{Buxton~RB, Uludag~K, Dubowitz~DJ, Liu~TT.} Modeling the hemodynamic response to brain activation. {NeuroImage}. 2004;23(Suppl.~1):S220--S233. \href{https://doi.org/10.1016/j.neuroimage.2004.07.013}{DOI:~10.1016/\allowbreak j.neuroimage.2004.07.013}.

	\bibitem{mandeville1999}
	{Mandeville~JB, Marota~JJA, Ayata~C, et al.} Evidence of a cerebrovascular postarteriole windkessel with delayed compliance. {J. Cereb. Blood Flow Metab}. 1999;19(6):679--689. \href{https://doi.org/10.1097/00004647-199906000-00012}{DOI:~10.1097/\allowbreak 00004647-\allowbreak 199906000-\allowbreak 00012}.

	\bibitem{friston1998b}
	{Friston~KJ, Fletcher~P, Josephs~O, et al.} Event-related fMRI: characterizing differential responses. {NeuroImage}. 1998;7(1):30--40. \href{https://doi.org/10.1006/nimg.1997.0306}{DOI:~10.1006/\allowbreak nimg.1997.0306}.

	\bibitem{friston1994}
	{Friston~KJ, Holmes~AP, Worsley~KJ, et al.} Statistical parametric maps in functional imaging: a general linear approach. {Hum. Brain Mapp}. 1994;2(4):189--210. \href{https://doi.org/10.1002/hbm.460020402}{DOI:~10.1002/\allowbreak hbm.460020402}.

	\bibitem{monti2011}
	{Monti~MM.} Statistical analysis of fMRI time-series: a critical review of the GLM approach. {Front. Hum. Neurosci}. 2011;5:28. \href{https://doi.org/10.3389/fnhum.2011.00028}{DOI:~10.3389/\allowbreak fnhum.2011.00028}.

	\bibitem{handwerker2012}
	{Handwerker~DA, Gonzalez-Castillo~J, D'Esposito~M, Bandettini~PA.} The continuing challenge of understanding and modeling hemodynamic variation in fMRI. {NeuroImage}. 2012;62(2):1017--1023. \href{https://doi.org/10.1016/j.neuroimage.2012.02.015}{DOI:~10.1016/\allowbreak j.neuroimage.2012.02.015}.

	\bibitem{lindquist2009}
	{Lindquist~MA, Meng Loh~J, Atlas~LY, Wager~TD.} Modeling the hemodynamic response function in fMRI: efficiency, bias and mis-modeling. {NeuroImage}. 2009;45(1, Suppl.~1):S187--S198. \href{https://doi.org/10.1016/j.neuroimage.2008.10.065}{DOI:~10.1016/\allowbreak j.neuroimage.2008.10.065}.

	\bibitem{polimeni2021}
	{Polimeni~JR, Lewis~LD.} Imaging faster neural dynamics with fast fMRI: a need for updated models of the hemodynamic response. {Prog. Neurobiol}. 2021;207:102174. \href{https://doi.org/10.1016/j.pneurobio.2021.102174}{DOI:~10.1016/\allowbreak j.pneurobio.2021.102174}.

	\bibitem{woolrich2004}
	{Woolrich~MW, Behrens~TEJ, Smith~SM.} Constrained linear basis sets for HRF modelling using Variational Bayes. {NeuroImage}. 2004;21(4):1748--1761. \href{https://doi.org/10.1016/j.neuroimage.2003.12.024}{DOI:~10.1016/\allowbreak j.neuroimage.2003.12.024}.

	\bibitem{wager2003}
	{Wager~TD, Nichols~TE.} Optimization of experimental design in fMRI: a general framework using a genetic algorithm. {NeuroImage}. 2003;18(2):293--309. \href{https://doi.org/10.1016/S1053-8119(02)00046-0}{DOI:~10.1016/\allowbreak S1053-\allowbreak 8119(02)00046-\allowbreak 0}.

	\bibitem{friston1998}
	{Friston~KJ, Josephs~O, Rees~G, Turner~R.} Nonlinear event-related responses in fMRI. {Magn. Reson. Med}. 1998;39(1):41--52. \href{https://doi.org/10.1002/mrm.1910390109}{DOI:~10.1002/\allowbreak mrm.1910390109}.

	\bibitem{robinson2006}
	{Robinson~PA, Drysdale~PM, Van der Merwe~H, et al.} BOLD responses to stimuli: dependence on frequency, stimulus form, amplitude, and repetition rate. {NeuroImage}. 2006;31(2):585--599. \href{https://doi.org/10.1016/j.neuroimage.2005.12.026}{DOI:~10.1016/\allowbreak j.neuroimage.2005.12.026}.

	\bibitem{birn2001}
	{Birn~RM, Saad~ZS, Bandettini~PA.} Spatial heterogeneity of the nonlinear hemodynamic response. {NeuroImage}. 2001;14(4):817--826. \href{https://doi.org/10.1006/nimg.2001.0873}{DOI:~10.1006/\allowbreak nimg.2001.0873}.

	\bibitem{mechelli2001}
	{Mechelli~A, Price~CJ, Friston~KJ.} Nonlinear coupling between evoked rCBF and BOLD signals: a simulation study of hemodynamic responses. {NeuroImage}. 2001;14(4):862--872. \href{https://doi.org/10.1006/nimg.2001.0876}{DOI:~10.1006/\allowbreak nimg.2001.0876}.

	\bibitem{magri2011}
	{Magri~C, Logothetis~NK, Panzeri~S.} Investigating static nonlinearities in neurovascular coupling. {Magn. Reson. Imaging}. 2011;29(10):1358--1364. \href{https://doi.org/10.1016/j.mri.2011.04.017}{DOI:~10.1016/\allowbreak j.mri.2011.04.017}.

	\bibitem{buxton1998}
	{Buxton~RB, Wong~EC, Frank~LR.} Dynamics of blood flow and oxygenation changes during brain activation: the Balloon model. {Magn. Reson. Med}. 1998;39(6):855--864. \href{https://doi.org/10.1002/mrm.1910390602}{DOI:~10.1002/\allowbreak mrm.1910390602}.

	\bibitem{friston2003}
	{Friston~KJ, Harrison~L, Penny~W.} Dynamic causal modelling. {NeuroImage}. 2003;19(4):1273--1302. \href{https://doi.org/10.1016/S1053-8119(03)00202-7}{DOI:~10.1016/\allowbreak S1053-\allowbreak 8119(03)00202-\allowbreak 7}.

	\bibitem{wager2005}
	{Wager~TD, Vazquez~A, Hernandez~L, Noll~DC.} Accounting for nonlinear BOLD effects in fMRI: parameter estimates and a model for prediction in rapid event-related studies. {NeuroImage}. 2005;25(1):206--218. \href{https://doi.org/10.1016/j.neuroimage.2004.11.008}{DOI:~10.1016/\allowbreak j.neuroimage.2004.11.008}.

	\bibitem{welvaert2014}
	{Welvaert~M, Rosseel~Y.} A review of fMRI simulation studies. {PLoS ONE}. 2014;9(7):e101953. \href{https://doi.org/10.1371/journal.pone.0101953}{DOI:~10.1371/\allowbreak journal.pone.0101953}.

	\bibitem{ellis2020}
	{Ellis~CT, Baldassano~C, Schapiro~AC, Cai~MB, Cohen~JD.} Facilitating open-science with realistic fMRI simulation: validation and application. {PeerJ}. 2020;8:e8564. \href{https://doi.org/10.7717/peerj.8564}{DOI:~10.7717/\allowbreak peerj.8564}.

	\bibitem{cakan2023}
	{Cakan~C, Jajcay~N, Obermayer~K.} neurolib: a simulation framework for whole-brain neural mass modeling. {Cogn. Comput}. 2023;15(4):1132--1152. \href{https://doi.org/10.1007/s12559-021-09931-9}{DOI:~10.1007/\allowbreak s12559-\allowbreak 021-\allowbreak 09931-\allowbreak 9}.

	\bibitem{havlicek2015}
	{Havlicek~M, Roebroeck~A, Friston~K, Gardumi~A, Ivanov~D, Uludag~K.} Physiologically informed dynamic causal modeling of fMRI data. {NeuroImage}. 2015;122:355--372. \href{https://doi.org/10.1016/j.neuroimage.2015.07.078}{DOI:~10.1016/\allowbreak j.neuroimage.2015.07.078}.

	\bibitem{sanzleon2015}
	{Sanz-Leon~P, Knock~SA, Spiegler~A, Jirsa~VK.} Mathematical framework for large-scale brain network modeling in The Virtual Brain. {NeuroImage}. 2015;111:385--430. \href{https://doi.org/10.1016/j.neuroimage.2015.01.002}{DOI:~10.1016/\allowbreak j.neuroimage.2015.01.002}.

	\bibitem{glover1999}
	{Glover~GH.} Deconvolution of impulse response in event-related BOLD fMRI. {NeuroImage}. 1999;9(4):416--429. \href{https://doi.org/10.1006/nimg.1998.0419}{DOI:~10.1006/\allowbreak nimg.1998.0419}.

	\bibitem{deneux2006}
	{Deneux~T, Faugeras~O.} Using nonlinear models in fMRI data analysis: model selection and activation detection. {NeuroImage}. 2006;32(4):1669--1689. \href{https://doi.org/10.1016/j.neuroimage.2006.03.006}{DOI:~10.1016/\allowbreak j.neuroimage.2006.03.006}.

	\bibitem{volterra1930}
	{Volterra~V.} Theory of Functionals and of Integral and Integro-Differential Equations. London: Blackie and Son; 1930.

	\bibitem{schetzen1980}
	{Schetzen~M.} The Volterra and Wiener Theories of Nonlinear Systems. New York: Wiley; 1980.

	\bibitem{marmarelis2004}
	{Marmarelis~VZ.} Nonlinear Dynamic Modeling of Physiological Systems. Hoboken: IEEE Press; Wiley-Interscience; 2004. \href{https://doi.org/10.1002/9780471679370}{DOI:~10.1002/\allowbreak 9780471679370}.

	\bibitem{hunter1986}
	{Hunter~IW, Korenberg~MJ.} The identification of nonlinear biological systems: Wiener and Hammerstein cascade models. {Biol. Cybern}. 1986;55(2--3):135--144. \href{https://doi.org/10.1007/BF00341929}{DOI:~10.1007/\allowbreak BF00341929}.

	\bibitem{boyd1985}
	{Boyd~S, Chua~LO.} Fading memory and the problem of approximating nonlinear operators with Volterra series. {IEEE Trans. Circuits Syst}. 1985;32(11):1150--1161. \href{https://doi.org/10.1109/TCS.1985.1085649}{DOI:~10.1109/\allowbreak TCS.1985.1085649}.

	\bibitem{pintelon2012}
	{Pintelon~R, Schoukens~J.} System Identification: A Frequency Domain Approach. 2nd ed. Hoboken: Wiley; 2012. \href{https://doi.org/10.1002/9781118287422}{DOI:~10.1002/\allowbreak 9781118287422}.

	\bibitem{schoukens2015}
	{Schoukens~J, Pintelon~R, Rolain~Y, et al.} Structure discrimination in block-oriented models using linear approximations: a theoretic framework. {Automatica}. 2015;53:225--234. \href{https://doi.org/10.1016/j.automatica.2014.12.045}{DOI:~10.1016/\allowbreak j.automatica.2014.12.045}.

\end{ltrtq}
}

\vspace{0.5\baselineskip}
\noindent{\bf Перевод подписей к рисункам / Figure captions}

\vspace{0.5\baselineskip}
\begin{enumerate}\setlength{\itemsep}{2pt}
\item \textbf{Fig.~1.} Linear impulse responses of the generators and the canonical SPM HRF ($C$), each normalized to its own maximum; the normalization is used only to compare the curve shapes. \textit{a} --- nominal P-DCM and BW. \textit{b} --- dependence of the response shape on parameters, shown for illustration: P-DCM with $\tau_v=0$ and with $\mu=0.8$, BW with $\mu=0$ and with the parameters from~\cite{friston2003}.
\item \textbf{Fig.~2.} Interaction of two events and deviation of the single-event response from the linear one for P-DCM and BW. \textit{a}, \textit{b} --- the response to a pair of events and the sum of the isolated responses at $a=0.30$ and SOA of 2~s; the difference between the curves equals $I(t)$, grey bars show the stimuli. \textit{c}, \textit{f} --- the interaction integral divided by $a^2$ versus SOA: the solid curve is the second order, the markers are the full computation at $a=0.05$, $0.15$, $0.30$. \textit{d}, \textit{g} --- $I(t)$ at $a=0.30$ and SOA of 2, 4 and 8~s, in percent of the baseline signal. \textit{e}, \textit{h} --- deviation of the single-event response from the linear one, divided by $a^2$, at the same amplitudes. In panels \textit{d}, \textit{e}, \textit{g}, \textit{h}, solid lines show the full computation and dashed lines the second order.
\item \textbf{Fig.~3.} Quadratic Volterra kernels $k_2(\tau_1,\tau_2)$ for P-DCM (\textit{a}--\textit{c}) and BW (\textit{d}--\textit{f}). The value at a point shows how the inputs $\tau_1$ and $\tau_2$ seconds ago jointly change the current response beyond the sum of their linear contributions: blue denotes attenuation, red enhancement; the colour scale is common. Columns: the full kernel (\textit{a}, \textit{d}), its static part, which depends only on the current value of the linear response (\textit{b}, \textit{e}), and the part it does not describe, $k_{2,\perp}$ (\textit{c}, \textit{f}). The diagonal corresponds to a single event; dashed lines in panels \textit{a}, \textit{d} correspond to event pairs with SOA of 2, 4 and 8~s (colours as in Fig.~2, \textit{d}, \textit{g}). The local second order is shown; the full finite-amplitude response also includes higher orders.
\item \textbf{Fig.~4.} Two regimes of nonlinear distortion of the BOLD signal in stimulus sequences: \textit{a} --- contrast underestimation at $a_A=0.30$, $a_B=0.05$ with random stimulus intervals and a mean SOA of 1.7~s; \textit{b} --- a nonzero contrast estimate at equal amplitudes $a_A=a_B=0.30$ and zero specified contrast, with condition A presented at short intervals and condition B at longer ones. The solid blue line is the P-DCM BOLD response, grey bars on the right axis show the stimulus times of conditions A and B, the green dashed line is the linear model with the canonical HRF at the input amplitudes, and the red dotted line is the standard GLM estimate; the panel labels give the coefficient and contrast estimates for the GLM and for the FV and FW corrections. Recording length 70~s, stimulus duration 1~s.
\item \textbf{Fig.~5.} Relative contrast error $100(\hat\Delta-\Delta)/\Delta$ for the intrinsic linear kernel $L_1$ and the canonical $C$ as a function of $a_A$ at $a_B=a_A/2$: P-DCM (\textit{a}--\textit{c}) and BW (\textit{d}--\textit{f}) at mean SOA of 2, 4 and 8~s. Markers are medians over three stimulus sequences, error bars show the min--max range between them.
\item \textbf{Fig.~6.} Dependence of the contrast error on the overall stimulation level and on the condition assignment for P-DCM and BW. \textit{a} --- error of $C$, $e=\hat\Delta-\Delta$, as the mean input amplitude $m=(a_A+a_B)/2$ grows at fixed $\Delta=0.10$; lines are medians over nine stimulus sequences, bands show the min--max range. \textit{b} --- error at $a_A=a_B=0.20$ and $\Delta=0$: assignment by the preceding interval with $C$ (diagonal crosses) and with $L_1$ (upright crosses), random assignment with $C$ (mean $\pm2$ standard deviations over 100 permutations); $k$ is the sequence number at a given mean SOA.
\item \textbf{Fig.~7.} Accuracy of the second order and the effect of substituting estimates into the corrections. \textit{a} --- contrast error of $C$ predicted by the second order versus the full computation for 90 combinations; the line denotes agreement, the grey band a discrepancy of at most 25\% or $0.001$ model units. \textit{b}, \textit{c} --- contrast errors $e=\hat\Delta-\Delta$ at $a_A=0.30$, $a_B=0.15$, the main sequence ($k=0$) and mean SOA of 2, 4 and 8~s for P-DCM (\textit{b}) and BW (\textit{c}). Bars show the canonical estimate $C$, the FV and FW corrections with the specified amplitudes (a theoretical control; in real data the amplitudes are unknown) and the same corrections with the $\hat\beta_C$ estimates substituted.
\end{enumerate}

\vspace{0.5\baselineskip}
\noindent{\bf Перевод подписей к таблицам / Table captions}

\vspace{0.5\baselineskip}
\noindent\textbf{Table~1.} Simulation runs addressing the research questions and the corresponding results.

\vspace{2pt}
\noindent\textbf{Table~2.} Relative contrast error for $C$ and the intrinsic $L_1$: median and min--max range over nine stimulus sequences; $a_B=a_A/2$.

\vspace{2pt}
\noindent\textbf{Table~3.} The FV correction in three generator regimes: pooled contrast RMSE for $C$ and FV and the relative change in RMSE, \%, by contrast class without noise (n/n) and at SNR 8 and 2. Negative values correspond to error reduction.

\vspace{2pt}
\noindent\textbf{Table~4.} The FV and FW corrections under the conditions of run 4b: relative change in contrast RMSE compared with $C$, \%, by contrast class without noise (n/n) and at SNR 8 and 2. Negative values correspond to error reduction.

\label{lastpage}
\end{document}
